\documentclass[10pt,a4paper]{amsart}
\usepackage[margin=19mm]{geometry}
\usepackage{amsmath,amssymb,mathtools}
\usepackage[scr=rsfs,cal=euler]{mathalfa}
\usepackage{xfrac}
\usepackage[unicode,hidelinks,bookmarks=false]{hyperref}
\newcommand{\ie}{i.e.\ }
\newcommand{\cf}{cf.\ }
\newcommand{\resp}{respectively\ }
\newcommand{\Subsection}{\subsection}
\providecommand{\nobreakdash}{\nobreak\hbox{-}}
\setlength{\emergencystretch}{2em}
\multlinegap0pt
\usepackage{tikz}
\newtheorem{prop}{Proposition}
\newcommand{\CF}{\mathcal{F}}
\newcommand{\SL}{\mathrm{SL}}
\newcommand{\Z}{\mathbb{Z}}
\newcommand{\hh}{\mathfrak{h}}
\title{Lattices, Gates, and Curves: GKP codes as a Rosetta stone}
\author{Jonathan Conrad, Ansgar G. Burchards, Steven T. Flammia}
\date{Source: arXiv:2407.03270v3 (CC BY 4.0)}
\begin{document}
\maketitle

\noindent\emph{Source and licence.} Lattices, Gates, and Curves: GKP codes as a Rosetta stone. Version arXiv:2407.03270v3 (CC BY 4.0), distributed by arXiv under the Creative Commons Attribution 4.0 International licence, http://creativecommons.org/licenses/by/4.0/, as recorded in the arXiv licence metadata for this exact version and shown on the abstract page https://arxiv.org/abs/2407.03270v3. Full licence notice: \texttt{licenses/arxiv-2407.03270-CC-BY-4.0.txt}.\par\smallskip

\noindent\textit{Editorial scope.} Counted unit: the article's prop prop:distance (source0.tex lines 633--644). The article's complete original proof of it is delivered with it (source0.tex lines 645--647, one \texttt{proof} environment of 3 lines). No mathematics is written, repaired or re-derived: the statement and the proof are the article's own lines, in the article's own notation, and no retained line is substituted or re-flowed.  No further source-local context is needed: every cross-reference of every retained line resolves inside the two delivered spans.  Nothing was omitted from any retained span, so the package carries no omission record.  No retained line contains a citation, so the wrapper carries no bibliography and no bibliography entry.  No retained line contains a figure environment, an image-inclusion command or drawing code, so no image is delivered and the package declares no asset dependency.  Editorial formatting only, in addition: an AMS \texttt{amsart} preamble that restates the article's own declarations and macros used by the retained lines, namely the environments \texttt{prop} on separate counters, together with the standard packages those lines need.  The retained environments print in this document as Proposition 1 in order of appearance, whereas the article prints the same items with the numbering of its own full document.  The complete native source of the article as distributed in the arXiv e-print for this exact version is bundled unchanged as \texttt{sources/161778/source0.tex}.
\begin{prop}[Distance and discriminant]\label{prop:distance}
For every nonempty set $T\subseteq\hh$,
\begin{equation}
\inf_{\tau\in T}\lambda_1(L_\tau)>0
\quad\Longleftrightarrow\quad
\inf_{\tau\in T}\|\Delta\|(\tau)>0.
\end{equation}
For the associated scaled single-mode GKP codes with fixed integer $d\geq2$, these conditions are also equivalent to
\begin{equation}
\inf_{\tau\in T}\Delta_{\rm GKP}(\tau)>0.
\end{equation}
\end{prop}

\begin{proof}
Both $\|\Delta\|$ and $\lambda_1(L_\tau)$ are $\SL_2(\Z)$-invariant, so we may represent the points of $T$ in the standard fundamental domain $\CF$. There the shortest vector of $\Z+\tau\Z$ has length $1$, and hence $\lambda_1(L_\tau)=y^{-1/2}$. Thus $\lambda_1$ is bounded away from zero on $T$ exactly when $y$ is bounded above. On every truncated fundamental domain $\CF\cap\{y\leq Y\}$, the continuous and nowhere-vanishing function $\|\Delta\|$ has a positive minimum. Conversely, the product expansion gives $\|\Delta\|(\tau)\leq C y^6e^{-2\pi y}\to0$ uniformly as $y\to\infty$. This proves the equivalence. The final statement follows from $\Delta_{\rm GKP}=d^{-1/2}\lambda_1$ for scaled codes.
\end{proof}

\end{document}
